\chapter{Main Results and Conclusion}
\begin{spacing}{1.15}

\section{Main Results}
Table~\ref{tab:main-results-summary} summarizes the main results of the
comparison study. A consistent pattern across all four models is that the
default threshold of 0.50 was too conservative for this intrusion-detection
task. After threshold tuning, every model achieved a tuned attack-class F1
above 0.90, which shows that threshold selection was an important part of the
final evaluation.

\begin{table}[htbp]
    \centering
    \small
    \setlength{\tabcolsep}{4pt}
    \caption{Main results, strengths, and weaknesses of the compared models.}
    \label{tab:main-results-summary}
    \begin{tabularx}{\linewidth}{>{\raggedright\arraybackslash}p{2.2cm} c c X X}
        \toprule
        Model & Test ROC-AUC & Tuned F1 & Main strength & Main weakness \\
        \midrule
        Logistic Regression & 0.9586 & 0.9080 & Most interpretable model; after tuning it reached 0.98 attack recall and remained competitive as a simple linear baseline. & Lowest ROC-AUC of the four models and the most sensitive to threshold choice. \\
        MLP & 0.9693 & 0.9121 & Strong nonlinear baseline; achieved 0.99 attack recall after tuning and very strong recall on Analysis (0.9529) and Fuzzers (0.9849). & Less interpretable than the linear model and still weaker than the tree ensembles on ROC-AUC. \\
        Random Forest & 0.9736 & 0.9091 & Strong tabular baseline with robust performance on most attack families and near-perfect recall on Backdoor, Worms, and Reconnaissance. & Weakest family-level recall on Analysis and Fuzzers (both 0.8614), which limited its final tuned F1. \\
        XGBoost & 0.9761 & 0.9154 & Best overall model; highest validation ROC-AUC, highest test ROC-AUC, and highest tuned F1, with strong results across nearly all attack families. & Most complex and least interpretable model, with a larger tuning burden than the baselines. \\
        \bottomrule
    \end{tabularx}
\end{table}

The results show a clear hierarchy. Logistic Regression served as a strong
interpretable baseline, but its lower ROC-AUC indicates that a linear decision
boundary cannot capture all of the structure in the UNSW-NB15 features. The
MLP improved substantially over the linear baseline and performed especially
well after threshold tuning, showing that a nonlinear neural model can compete
on this task. However, the two tree-based ensemble models achieved the
strongest threshold-independent ranking performance, which suggests that the
dataset contains important nonlinear feature interactions that are well handled
by ensemble trees.

Per-attack-category analysis also reveals where the models struggle. Analysis
was the hardest attack family for all four models: recall was 0.9586 for
Logistic Regression, 0.9529 for MLP, 0.8614 for Random Forest, and 0.8645 for
XGBoost. Fuzzers was another challenging category for some models: Logistic
Regression reached 0.9591 recall, Random Forest 0.8614, MLP 0.9849, and
XGBoost 0.9648. In contrast, Generic, Backdoor, Worms, and Exploits were
detected extremely well by the stronger models, usually with recall close to
1.0.

\section{Conclusion}
Our final selected model is \textbf{XGBoost}. It achieved the best
cross-validated ROC-AUC (0.9897), the best test ROC-AUC (0.9761), and the best
tuned F1 (0.9154), which makes it the strongest overall solution to the binary
network intrusion detection problem studied in this project. Random Forest was
the closest competitor and remains a strong alternative when simpler ensemble
behavior is preferred. Logistic Regression remains valuable as an interpretable
baseline, and MLP demonstrates that neural models can be competitive on tabular
intrusion data when tuned carefully. Overall, the study shows that while all
four models can achieve strong tuned performance, boosted tree ensembles offer
the best balance of detection quality, robustness across attack families, and
practical predictive performance on UNSW-NB15.

\end{spacing}
\newpage
